%% =====================================================================
%%  Corrigés — Unit 4 « This is my family »
%% =====================================================================
\corrigesUnite{4}{This is my family}

\begin{corrige}{1}
\begin{reponses}[4]
\item \rep{father}
\item \rep{brother}
\item \rep{son}
\item \rep{husband}
\item \rep{uncle}
\item \rep{nephew}
\item \rep{grandfather}
\item \rep{boy}
\end{reponses}
\retiens{\en{grand-} se place devant \en{mother}, \en{father} et \en{parents} : \en{grandmother},
\en{grandfather}, \en{grandparents}.}
\end{corrige}

\begin{corrige}{2}
\begin{center}\otArbreMartin[solution]\end{center}
\begin{reponses}[4]
\item \rep{mother}
\item \rep{brother}
\item \rep{sister}
\item \rep{husband}
\item \rep{son}
\item \rep{aunt}
\item \rep{grandfather}
\item \rep{uncle}
\end{reponses}
\pourquoi{lisez chaque phrase « à l'envers » : \en{Léa is Nathan's aunt} = Léa est la tante de
Nathan. Le possesseur (\en{Nathan's}) est la personne de référence ; le mot qui suit dit ce qu'est
l'autre personne pour elle. Les liens sont notés sous chaque prénom, par rapport à Léa.}
\end{corrige}

\begin{corrige}{3}
\begin{reponses}[4]
\item \rep{my}
\item \rep{your}
\item \rep{his}
\item \rep{her}
\item \rep{Its}
\item \rep{our}
\item \rep{their}
\item \rep{her}
\end{reponses}
\pourquoi{Mr.\ Carter est un homme → \en{his} ; Emma et Léa sont des femmes → \en{her} ; le chien
→ \en{its} ; Sofia et David (deux personnes) → \en{their}.}
\piege{écrire \en{It's name}. \en{it's} signifie toujours \en{it is} ; le possessif \en{its} n'a
pas d'apostrophe.}
\end{corrige}

\begin{corrige}{4}
\begin{reponses}[2]
\item \rep{Léa's brother}
\item \rep{Jack's sister}
\item \rep{Sofia's husband}
\item \rep{my brother's dog}
\item \rep{the child's name}
\item \rep{my parents' house}
\end{reponses}
\pourquoi{pas d'article devant le possesseur quand c'est un prénom (\en{Léa's}, pas \en{the
Léa's}). \en{parents} se termine déjà par \en{-s} : on ajoute seulement l'apostrophe.}
\end{corrige}

\begin{corrige}{5}
\begin{reponses}[1]
\item \rep{Léa and her brother are in Lyon.}
\item \rep{Jack is with his sister.}
\item \rep{This is my brother's dog.}
\item \rep{Their names are Leo and Max.}
\item \rep{The cat? Its name is Mango.}
\item \rep{Emma's parents are British.}
\end{reponses}
\pourquoi{phrases 4 et 6 : le sujet est au pluriel (\en{names}, \en{parents}), donc \en{are}.}
\end{corrige}

\begin{corrige}{6}
\begin{reponses}[3][french]
\item \rep{F}
\item \rep{V}
\item \rep{V}
\item \rep{F}
\item \rep{V}
\item \rep{F}
\end{reponses}
\pourquoi{la famille d'Emma est à Londres ; la femme de Tom est irlandaise (\en{Irish}) ; Emma est
célibataire (\en{I'm single}).}
\end{corrige}

\begin{corrige}{7}
\begin{reponses}[4]
\item \rep{they}
\item \rep{den}
\item \rep{there}
\item \rep{dough}
\item \rep{those}
\item \rep{dares}
\item \rep{breeze}
\item \rep{worthy}
\end{reponses}
\pourquoi{le son \phon{ð} est plus « doux » et plus long que \phon{d} : on l'entend « frotter ».
Texte de la piste page~\pageref{trans-4.10}.}
\end{corrige}

\begin{corrige}{8}
\begin{reponses}[4]
\item \rep{mother}
\item \rep{58}
\item \rep{father}
\item \rep{61}
\item \rep{sister}
\item \rep{27}
\item \rep{nephew}
\item \rep{3}
\end{reponses}
\pourquoi{Josh est le fils de Kate (\en{her son}), la sœur de Jack : c'est donc son neveu, comme
Jack le dit lui-même : \en{He's my nephew!}}
\end{corrige}

\begin{corrige}{9}
\begin{anglais}\otTableau
\begin{tabular}{@{}l@{\hspace{6mm}}l@{}}
\otNomDit{lea} & Sam, is that \rep{your} sister?\\
\otNomDit{sam} & Yes, it is. \rep{Her} name is Maya.\\
\otNomDit{lea} & How \rep{old} is she?\\
\otNomDit{sam} & She's \rep{twenty}. And this is my \rep{brother}, Ben.\\
\otNomDit{lea} & Are your \rep{parents} in Brooklyn too?\\
\otNomDit{sam} & No, they \rep{aren't}. They're in Atlanta.\\
\otNomDit{lea} & And who's \rep{this}?\\
\otNomDit{sam} & That's \rep{our} cat, Mango!\\
\end{tabular}
\end{anglais}
\retiens{\en{our cat} : le chat appartient à toute la famille de Sam (« notre chat »).}
\end{corrige}

\begin{corrige}{10}
\begin{reponses}[1]
\item \rep{This is my mother and my father.} \textit{ou} \rep{These are my mother and father.}
\item \rep{Emma and her brother are British.}
\item \rep{Jack and his sister are in New York.}
\item \rep{Sofia's husband is Canadian.}
\item \rep{Their children are ten and seven.}
\item \rep{What a nice photo! Who's this?}
\item \rep{It's my parents' cat.} \textit{ou} \rep{That's my parents' cat.}
\item \rep{I'm an only child.}
\end{reponses}
\pourquoi{phrase 1 : \en{These are…} (« Voici… », au pluriel) sera étudié à l'unité 5. Phrase 5 :
l'âge s'exprime avec \en{be}, au pluriel \en{are}.}
\piege{oublier l'article dans « enfant unique » : on dit \en{an only child}.}
\end{corrige}

\begin{corrige}{11}
Réponse modèle (la famille de Jack) :
\begin{anglais}
\begin{tcolorbox}[otBase,colback=otGreenTint,colframe=otGreenTint,fontupper=\fontsize{10.4}{14}\selectfont]
This is my family. My mother's name is Linda. She's fifty-eight. My father, Bob, is sixty-one.
They're American and they're from New York. My sister's name is Kate. She's twenty-seven and she's
married. Her husband's name is Ryan. Their son, Josh, is three. He's so cute!
\end{tcolorbox}
\end{anglais}
\retiens{variez les tournures : \en{This is…}, \en{My mother's name is…}, \en{Her husband's name
is…}, et n'oubliez pas \en{be} pour l'âge.}
\end{corrige}

\begin{corrige}{12}
Réponses modèles :
\begin{reponses}[1]
\item \rep{Yes, it is. It's in Lyon.} \textit{ou} \rep{No, it isn't. It's in Belgium.}
\item \rep{Her name is Claire.} \textit{ou} \rep{My mother's name is Claire.}
\item \rep{His name is Paul.}
\item \rep{No, I'm not. I'm single.} \textit{ou} \rep{Yes, I am.}
\item \rep{Thank you!}
\end{reponses}
\retiens{pour parler de votre mère, \en{her name} ; de votre père, \en{his name}.}
\end{corrige}

\begin{corrigetest}
\begin{reponses}[5][french]
\item \rep{a}
\item \rep{b}
\item \rep{b}
\item \rep{b}
\item \rep{c}
\item \rep{a}
\item \rep{b}
\item \rep{a}
\item \rep{b}
\item \rep{a}
\end{reponses}
\pourquoi{questions 2 et 3 : \en{his} ou \en{her} selon le possesseur (Léa → \en{her}, Jack →
\en{his}). Question 5 : \en{their} est invariable et \en{child} a un pluriel irrégulier,
\en{children}.}
\end{corrigetest}
